\section{New local checks in Masullo's Chapters 11--12}
\label{sec:local-source-audit}

\subsection{Scope and source normalization}

This section checks the public 940-page draft of Masullo, using printed
page numbers, which agree with its PDF page numbers.  The statements and
proofs on pp.~491--496 and 514--519 were read, together with the relevant
definitions on pp.~243, 259--260 and 427--428, Theorem~7.4.2.1 on
p.~367, and Lemma~C.3.2.2 on pp.~695--696.  The displayed formulas on
pp.~496, 514 and 517 were also checked in rendered page images.  No
copyrighted source PDF or page image is included in this deliverable.
These checks do not assume either $A$ or $M$, or any assertion concerning
the location of nontrivial zeta zeros.  They are additional to the
previously identified indentation and global-to-local problems.

The conclusions have different scopes.  MA-LA-001 is a literal
normalization inconsistency in a Poisson-summation step.  MA-LA-002 is
a false codomain for the printed cutoff.  MA-LA-003--004 establish a
path dependence for the natural annular repair, and give a sufficient
rate that repairs this \emph{local} trace limit.  MA-LA-005 identifies
two false intermediate kernel equalities but also explains their
cancellation in the real, even calculation; it is not an additional
counterexample to the repaired limiting trace formula.  MA-LA-006
separates the valid cyclic nuclear-trace lemma from an unverified
factorization hypothesis in its application.

\subsection{The missing half-weight in the orbital Poisson sum}

\paragraph{MA-LA-001 (the centered and uncentered orbital sums differ).}
For $h\in C_c^\infty(\mathbb R_{>0})$, set
\[
 \mathcal Mh(s)=\int_0^\infty h(t)t^s\,\frac{dt}{t}.
\]
For every rational prime $p$,
\[
 \sum_{n\in\mathbb Z}
 \mathcal Mh\!\left(\frac{2\pi i n}{\log p}-\frac12\right)
 =\log p\sum_{k\in\mathbb Z}p^{-k/2}h(p^k).
 \tag{MA-LA.1}
\]
Both sides are absolutely convergent.  In particular, if $h$ is
supported in $(3/2,5/2)$ and $h(2)=1$, then at $p=2$ the value is
$\log 2/\sqrt2$, whereas the unweighted sum is $\log 2$.

\emph{Proof.}
Put $L=\log p$ and $H(u)=e^{-u/2}h(e^u)$.  This is a compactly
supported smooth function on $\mathbb R$.  For the Fourier convention
$\widehat H(\xi)=\int H(u)e^{-2\pi i u\xi}\,du$, the summand on
the left is $\widehat H(-n/L)$.  Poisson summation gives
$\sum_n\widehat H(n/L)=L\sum_k H(kL)$, proving (MA-LA.1).
Rapid decay of $\widehat H$ proves absolute convergence; the other
sum is finite.  The support condition leaves only $2^1$.
\hfill$\square$

Here is the precise comparison with the draft.  In Theorem~7.4.2.1,
equation~(7.4.2.1), p.~367, $\varepsilon_p$ is defined by the zeros of
the \emph{ordinary} reciprocal local Euler factor, and the integrand
has $|g|^{\varepsilon_p}$.  For $F=\mathbb Q$ and $\chi=1$ these
zeros are $2\pi i n/\log p$, each of multiplicity one.  In the overview
on p.~494 and equation~(11.3.2.5), p.~496, the same test function and
the same local-zero parameters occur with
$|g|^{\varepsilon_p-1/2}$.  The following sentence identifies this
with the earlier right side, although its exponent differs by
$-1/2$.  Equation~(11.3.2.1), p.~495, and the Poisson conclusion on
p.~494 give an unweighted prime-power sum.

Take $\varphi(g)=h(|g|)$, normalize the compact norm-one fiber to
volume one, and use the standard unramified measure
$\mu_p(\mathbb Z_p^\times)=1$.  Then the displayed centered sum in
(11.3.2.5) is exactly the left side of (MA-LA.1).  A local uniformizer
has norm $p^{-1}$, so the printed unweighted sum
$\log p\sum_k\varphi(\pi_p^k)$ is
$\log p\sum_k h(p^k)$ after reindexing.  The corrected centered sum
in uniformizer notation is
\[
 \log p\sum_{k\in\mathbb Z}
 |\pi_p^k|^{-1/2}\varphi(\pi_p^k)
 =\log p\sum_{k\in\mathbb Z}p^{-k/2}h(p^k).
 \tag{MA-LA.2}
\]
Thus no implicit twist found in these source definitions removes the
discrepancy: twisting the test changes its value on the printed right
side as well.  Deleting the centering or inserting the half-weight are
different edits.  Inserting the weight repairs this Poisson step; it
does not repair the indentation or establish a global-to-local
intertwiner.  The example uses only zeros of $1-p^{-s}$ and Poisson
summation, not zeros of a global zeta function.

\subsection{The cutoff at the additive origin}

\paragraph{MA-LA-002 (an outer cutoff does not have the stated codomain).}
The source defines
\[
 \mathcal S_{00}(\mathbb R)
 =\{f\in\mathcal S(\mathbb R):f^{(j)}(0)=0\text{ for all }j\ge0\}
\]
in Theorem~5.3.2.1(1), p.~243; Definition~5.3.4.1, p.~259,
incorporates this factor at real places.  A smooth multiplier $a$ with
the usual Schwartz multiplier bounds satisfies
\[
 M_a\mathcal S(\mathbb R)\subseteq\mathcal S_{00}(\mathbb R)
 \quad\Longleftrightarrow\quad
 a^{(j)}(0)=0\quad(j\ge0).
 \tag{MA-LA.3}
\]
Consequently an outer cutoff equal to one near zero does not define
the map $P_{\rho,\lambda}:\mathcal S\to\mathcal S_{00}$ asserted
in Step II on p.~516.

\emph{Proof.}
Necessity follows by applying the multiplier to a test function equal
to one in a neighborhood of zero.  Sufficiency follows from the
Leibniz rule.  For the printed cutoff $a(0)=1$, even the zeroth
condition fails.\hfill$\square$

The preceding Fourier cutoff does not remove this issue, even on the
initial polar-free, zero-integral domain.  Let $b$ be a nonzero, even,
nonnegative compactly supported smooth frequency cutoff, and put
$q=\mathcal F^{-1}b$.  The function $q$ is nonconstant and even.
Choose two even smooth bumps supported away from zero, each of
integral one, whose integrals against $q$ differ; this is possible
by concentrating them near two positive points where $q$ differs.
Their difference $f$ belongs to $\mathcal S_{00}$ and has integral
zero, but
\[
 (\mathcal F^{-1}M_b\mathcal Ff)(0)=\int q(-y)f(y)\,dy\ne0.
\]
If $h\ge0$ is nonzero and supported in $(1,\infty)$, subsequent
averaging by $t^{-1/2}f(x/t)$ multiplies this value at zero by the
positive number $\int h(t)t^{-1/2}\,dt/t$.  An outer spatial cutoff
preserves the nonzero value.  Thus merely starting the composition
on $\mathcal S_{00}$ does not justify its claimed final codomain.

The footnote on p.~514 explicitly proposes replacing the spatial
cutoff by a smooth annulus $[1/\lambda'',\lambda]$, with
$\lambda''$ independent of $\lambda$, and asks for a convergence
study.  Vanishing on a neighborhood of zero does repair (MA-LA.3).
The next two results supply the required convergence study in the
real, trivial-isotype case and show why the relative rate matters.

\subsection{Exact real-place trace and the annular limit}

We use the Fourier normalization
\[
 (\mathcal F f)(\xi)=\int_{\mathbb R}f(x)e^{-2\pi i x\xi}\,dx,
 \qquad \check B=\mathcal F^{-1}B.
\]
Let $h\in C_c^\infty((1,\infty))$.  For $t>0$ set
\[
 D_tf(x)=t^{-1/2}f(x/t),\qquad
 A_h=\int_0^\infty h(t)D_t\,\frac{dt}{t}.
\]
Let $A,B\in C_c^\infty(\mathbb R)$ be even, satisfy $0\le A,B\le1$,
and equal one on a neighborhood of zero.  For $X,Y,\epsilon>0$, put
\[
 a_{X,\epsilon}(x)=A(x/X)\bigl(1-A(x/\epsilon)\bigr),
 \qquad b_Y(\xi)=B(\xi/Y),
 \qquad T_{X,Y,\epsilon}=M_{a_{X,\epsilon}}A_h
       \mathcal F^{-1}M_{b_Y}\mathcal F.
 \tag{MA-LA.4}
\]

\paragraph{MA-LA-003 (exact trace and two distinct annular limits).}
The operator in (MA-LA.4) is trace class, commutes with reflection,
and on $L^2_{\mathrm{even}}(\mathbb R,dx)$ has trace
\begin{align*}
 \operatorname{Tr}_{\mathrm{even}}T_{X,Y,\epsilon}
  =\frac12\int h(t)t^{-1/2}
   \int a_{X,\epsilon}(x)Y
   \left\{\check B\bigl(Y(t^{-1}-1)x\bigr)
          +\check B\bigl(Y(t^{-1}+1)x\bigr)\right\}
      dx\,\frac{dt}{t}.
 \tag{MA-LA.5}
\end{align*}
Write $L=XY$ and $r=Y\epsilon$.  If $L\to\infty$, then
\[
 \operatorname{Tr}_{\mathrm{even}}T_{X,Y,\epsilon}
 \longrightarrow
 \begin{cases}
 A_0(h):=\displaystyle\int_1^\infty
             h(t)\frac{\sqrt t}{t^2-1}\,dt,&r\to0,\\[6pt]
 0,&r\to\infty.
 \end{cases}
 \tag{MA-LA.6}
\]
For fixed $A,B,h$, the first error is
$O_N(L^{-N})+O(r)$, and the second is $O_N(r^{-N})$, for every
fixed $N$, after increasing the constants.  In particular the first
limit is strictly positive for nonnegative nonzero $h$.

\emph{Proof.}
The kernel, with respect to additive Lebesgue measure, is
\[
 K_T(x,y)=a_{X,\epsilon}(x)
 \int h(t)t^{-1/2}Y\check B\bigl(Y(x/t-y)\bigr)\,\frac{dt}{t}.
 \tag{MA-LA.7}
\]
It is a Schwartz function of $(x,y)$: $x$ has compact support and,
uniformly for $x$ there and $t$ in the fixed compact support of $h$,
all derivatives decay rapidly as $|y|\to\infty$.  A Schwartz kernel
is trace class and has trace its diagonal integral.  For completeness,
expand it in Hermite functions.  Applying arbitrary powers of the
harmonic oscillator in each variable gives coefficients
$c_{mn}=O_N((1+m)^{-N}(1+n)^{-N})$.  Their absolute sum is finite,
so the rank-one expansion converges in trace norm.  Moreover
$\int|e_m(x)e_n(x)|\,dx\le1$, which justifies integrating the
diagonal expansion and identifies its trace.

Let $Jf(x)=f(-x)$.  Evenness of both cutoffs makes $T$ commute with
$J$.  The even-sector trace is
$\tfrac12(\operatorname{Tr}T+\operatorname{Tr}(TJ))$; the two
diagonals are $K_T(x,x)$ and $K_T(x,-x)$.  This proves (MA-LA.5).

Substitute $z=Yx$.  Its inner integral becomes
\[
 \int A(z/L)(1-A(z/r))
 \{\check B(c_-(t)z)+\check B(c_+(t)z)\}\,dz,
 \qquad c_\pm(t)=t^{-1}\pm1.
 \tag{MA-LA.8}
\]
Because the support of $h$ is a compact subset of $(1,\infty)$,
both $|c_-|$ and $|c_+|$ have a positive lower bound there.  Thus
the Schwartz functions in (MA-LA.8) have an integrable dominating
function uniform in $t$.  If $r\to0$, the multiplier tends to one
for $z\ne0$; if $r\to\infty$, it tends to zero.  Dominated
convergence applies.  In the first case Fourier inversion and
$B(0)=1$ give
\[
 \int\check B(c_\pm(t)z)\,dz=\frac1{|c_\pm(t)|}.
\]
Consequently the limiting trace is
\[
 \frac12\int h(t)t^{-1/2}
 \left(\frac{t}{t-1}+\frac{t}{t+1}\right)\frac{dt}{t}
 = A_0(h).
\]
The outer-cutoff error is bounded by a Schwartz tail beyond a
constant times $L$.  For small $r$, the removed inner part is
supported on an interval of length $O(r)$ and its integrand is
uniformly bounded.  For large $r$, the remaining integrand vanishes
on an interval of radius a constant times $r$ and is bounded by a
Schwartz tail.  These prove the error estimates.\hfill$\square$

This is the normalized kernel that the draft itself writes in
(12.4.1.5) and (12.4.1.7), p.~516, followed by the three operators
listed in Step II and composed on p.~517.  Indeed, for
$F=\mathbb Q$, $S=\{\infty\}$, and the trivial compact isotype,
$\mathcal O_{\mathbb Q,S}^\times=\{\pm1\}$ and its invariant
sector is the even sector.  The normalized dilation in those
equations is precisely $t^{-1/2}f(x/t)$.  One may choose the lift
$\widetilde\varphi(t)=h(t)$ for $t>0$ and zero for $t<0$;
periodization by $\{\pm1\}$ gives $\varphi(g)=h(|g|)$.
Summing the two diagonals over $x>0$ equals the half of the full
real integral in (MA-LA.5).  The right side of the source's
(12.4.1.1), specialized to this case, is
\[
 \frac12\int_{\mathbb R^\times}
 h(|g|)\frac{|g|^{1/2}}{|1-g|}\,\frac{dg}{|g|}=A_0(h).
 \tag{MA-LA.9}
\]
There is no principal-value issue or logarithmic term for this $h$,
since its support is separated from $1$.  Thus the calculation tests
the actual displayed local kernel, without assuming the earlier
identification of a global geometric operator with that kernel.

\paragraph{MA-LA-004 (the path dependence persists for the source's
thin boundary smoothings).}
The fixed profiles in (MA-LA.4) are not needed for the obstruction.
Let $a_{X,\epsilon,\delta}$ be an even smooth approximation to
$1_{\{\epsilon\le|x|\le X\}}$, and $b_{Y,\delta}$ an even smooth
approximation to $1_{[-Y,Y]}$, with values in $[0,1]$, differing
from the indicators only in boundary intervals of total length
$O(\delta)$.  They may be chosen supported in the indicated
annulus and interval and equal to one outside boundary layers of
width $O(\delta)$.  Suppose
\[
 \delta<\tfrac14\epsilon,\qquad
 \delta(X+Y)\longrightarrow0.
 \tag{MA-LA.10}
\]
Then the trace of the same composition with these cutoffs differs
by $O_h(\delta(X+Y))$ from
\[
 \frac12\int h(t)t^{-1/2}
 \sum_{c\in\{c_-(t),c_+(t)\}}
 \frac{2}{\pi|c|}
 \left\{\operatorname{Si}(2\pi|c|XY)
       -\operatorname{Si}(2\pi|c|Y\epsilon)\right\}\,\frac{dt}{t},
 \tag{MA-LA.11}
\]
where $\operatorname{Si}(u)=\int_0^u(\sin v)/v\,dv$.
In particular, both paths
\[
 X=Y=N,\quad\delta=N^{-4},\qquad
 \epsilon=N^{-2}\quad\hbox{or}\quad\epsilon=N^{-1/2}
 \tag{MA-LA.12}
\]
satisfy the thin-smoothing requirements, but their traces tend to
$A_0(h)$ and $0$, respectively.

\emph{Proof.}
The inverse Fourier transform of the sharp frequency interval is
\[
 q_Y(v)=\frac{\sin(2\pi Yv)}{\pi v},\qquad q_Y(0)=2Y.
\]
For $c\ne0$, its integral over the sharp spatial annulus is
\[
 \int_{\epsilon\le|x|\le X}q_Y(cx)\,dx
 =\frac{2}{\pi|c|}
  \{\operatorname{Si}(2\pi|c|YX)
        -\operatorname{Si}(2\pi|c|Y\epsilon)\}.
\]
We use this merely as an explicit comparison integral; no trace
class assertion for discontinuous cutoff operators is required.
The frequency smoothing changes its inverse Fourier transform
uniformly by at most
$\|b_{Y,\delta}-1_{[-Y,Y]}\|_1=O(\delta)$, and integration over
the spatial support costs $O(X)$.
The spatial smoothing changes its $L^1$ norm by $O(\delta)$,
while $\|\mathcal F^{-1}b_{Y,\delta}\|_\infty\le
\|b_{Y,\delta}\|_1=O(Y)$.  Apply these estimates to each summand
in (MA-LA.5); the $t$ factors are bounded on the support of $h$.
This proves the stated error.

The classical improper integral
$\operatorname{Si}(u)\to\pi/2$ as $u\to\infty$, with error
$O(u^{-1})$ by integration by parts.  Also
$\operatorname{Si}(u)=O(u)$ as $u\to0$.  The constants are uniform
for the present $c_\pm(t)$.  Therefore (MA-LA.11) tends to
$A_0(h)$ when $XY\to\infty$, $Y\epsilon\to0$, and to zero
when $XY\to\infty$, $Y\epsilon\to\infty$.
All smoothing errors in (MA-LA.12) tend to zero.\hfill$\square$

Writing $X=\lambda$, $Y=\lambda'$ and
$\epsilon=1/\lambda''$, the essential ratio is
\[
 r=\frac{\lambda'}{\lambda''}.
 \tag{MA-LA.13}
\]
Both paths in (MA-LA.12) have
$\lambda,\lambda',\lambda''\to\infty$ and
$\delta\lambda,\delta\lambda'\to0$, the latter matching the
thin-smoothing regimes stated at the end of the source calculation
on p.~519.  Thus simultaneous divergence, even with those
smoothing conditions, is insufficient for the annular repair.
The sufficient rate $\lambda'/\lambda''\to0$ does restore the
local formula for tests supported away from $1$.  For the thin
boundary smoothing above the error is explicitly
\[
 O_h((\lambda\lambda')^{-1})+
 O_h(\lambda'/\lambda'')+
 O_h(\delta(\lambda+\lambda'))
 \tag{MA-LA.14}
\]
in this regime.  This result does not assert necessity of that rate
for each individual sign-changing test, and it does not settle
tests meeting $1$, the subtraction there, general $S$, or the
global intertwining step.

There is also a direct reading of the source's proposed volume
correction.  With a sharp annular spatial cutoff, the allowed
positive coordinate after $w=xu$ is
\[
 \max\{|w|/\lambda',1/\lambda''\}\le x\le\lambda.
\]
Its multiplicative volume is
\[
 \left[\log\lambda-
   \log\max\{|w|/\lambda',1/\lambda''\}\right]_+
 =\left[\log(\lambda\lambda')-
   \log\max\{|w|,r\}\right]_+.
\]
This is the correction suggested in the footnote on p.~514 to
(12.4.1.23), p.~518.  It displays the same residual parameter $r$;
one cannot discard it merely because both cutoffs diverge.

\subsection{The transposed kernel and why its errors can cancel}

\paragraph{MA-LA-005 (false intermediate equalities, with a limited
cancellation).}
The kernel convention in (12.4.1.5)--(12.4.1.7) is
\[
 K(x,y)=\int h(t)t^{-1/2}\delta(y-t^{-1}x)\,\frac{dt}{t}.
 \tag{MA-LA.15}
\]
For fixed $x$ and $\xi\in\{\pm1\}$, its substitution into
$K(x+z,\xi x)$ gives
\[
 \delta(\xi x-t^{-1}(x+z))
   =t\,\delta(z-(t\xi-1)x).
 \tag{MA-LA.16}
\]
Equation~(12.4.1.13), p.~517, instead substitutes
$\delta(x+z-t^{-1}\xi x)$, the kernel with the two arguments
exchanged.  Its factor is $t^{-1/2}$, whereas the actual integral
of the left kernel against a Fourier character has factor
$t^{1/2}$ and phase $(t\xi-1)xu$.

\emph{Verification.}
Equation (MA-LA.16) is the one-dimensional scaling rule for the
Dirac distribution, $\delta(t^{-1}v)=t\delta(v)$ for $t>0$.
Thus, with the negative Fourier character used in the source,
\[
 \int K(x+z,\xi x)e^{-2\pi i zu}\,dz
  =\int h(t)t^{1/2}
          e^{-2\pi i(t\xi-1)xu}\,\frac{dt}{t}.
 \tag{MA-LA.17}
\]
This is not its printed right side.\hfill$\square$

There is a preceding transposition as well.  The correct kernel
composition on p.~517 is
$a(x)\int K(x,z)q(z-y)\,dz$.  Its diagonal at $y=\xi x$
therefore uses $K(x,\xi x+z)$ against $q(z)$, not
$K(x+z,\xi x)$ against $q(z)$.  Parseval in its bilinear form
requires opposite Fourier signs.  The source's two characters in
(12.4.1.12) both have negative signs; that sign issue disappears
here because $b$, and hence $q$, is even.  Evenness does not by
itself justify interchanging the arguments of $K$.

It is necessary to record the cancellation, rather than count two
typographical slips as a further failure of the final real local
formula.  After the erroneous delta substitution, the source uses
the phase $(t^{-1}\xi-1)xu$ with the correct coefficient
$t^{-1/2}$.  For $\xi=1$ this is precisely the correct diagonal
phase.  For $\xi=-1$ its sign is the negative of the correct phase,
which again has no effect against even $b$.  Thus replacing the
two false intermediate lines by the direct kernel calculation
(MA-LA.7) yields (MA-LA.5).  The later variable change $t\mapsto
t^{-1}$ is legitimate with multiplicative Haar measure, and the
limiting density $|t|^{1/2}/|1-t|$ is inversion invariant.  These
facts explain why the source can arrive at the correct local
density despite the intermediate transpositions.  They do not
remove the annular path dependence of MA-LA-003--004.

The first erroneous transposition is not an exact finite-cutoff
identity that happens to be disguised.  For example, under the
fixed-profile scaling of MA-LA-003, with $L\to\infty$ and
$r\downarrow0$, its even trace would replace $t^{-1/2}$ and
$t^{-1}\pm1$ by $t^{1/2}$ and $t\pm1$.  Both expressions have
the same limit at $r=0$, but the coefficients of the removed inner
interval, to first order in $r$, are respectively
\[
 -r\check B(0)\!\left(\int A\right)
       \int h(t)t^{-1/2}\,\frac{dt}{t},
 \qquad
 -r\check B(0)\!\left(\int A\right)
       \int h(t)t^{1/2}\,\frac{dt}{t}.
\]
They differ for $h\ge0$, $h\ne0$ supported above $1$.
Indeed, substitute $z=rv$ into the removed term and use continuity
of $\check B$ at zero; $\check B(0)=\int B>0$.  Hence the
intermediate equality fails already for smooth local cutoffs.

\subsection{What is required for the trace factorization on p.~515}

\paragraph{MA-LA-006 (the cyclic trace lemma is valid; range inclusion
alone is insufficient).}
Lemma~C.3.2.2, pp.~695--696, requires a \emph{nuclear} lift
$A_0:E_2\to E_1$ through a continuous injection $T:E_1\to E_2$,
not only a nuclear endomorphism $A=TA_0$ of $E_2$ whose range
lies in $T(E_1)$.  Under its hypotheses, the claimed equality
\[
 \operatorname{Tr}(A_0T)=\operatorname{Tr}(TA_0)
 \tag{MA-LA.18}
\]
has a direct proof.

\emph{Proof.}
Choose a nuclear presentation
$A_0x=\sum_n\lambda_n\phi_n(x)y_n$ with
$\sum|\lambda_n|<\infty$, $(\phi_n)$ equicontinuous and
$(y_n)$ bounded.  The two compositions have presentations
\[
 TA_0x=\sum_n\lambda_n\phi_n(x)Ty_n,
 \qquad
 A_0Tz=\sum_n\lambda_n(\phi_n\circ T)(z)y_n.
\]
Their traces, defined independently of presentation using the
approximation properties assumed in the lemma, are the identical
absolutely convergent series
$\sum_n\lambda_n\phi_n(Ty_n)$.
\hfill$\square$

The finite-rank approximation explanation on p.~696 invokes
convergence uniform on precompact sets.  This topology by itself
does not make trace continuous: on $\ell^2$,
$n^{-1}P_{\{1,\ldots,n\}}\to0$ even in operator norm, but each
trace is one.  The nuclear presentation above supplies the needed
stronger control and proves the lemma without that inference.
Accordingly this observation is a repair of a proof detail, not a
counterexample to the lemma.

Nor does nuclearity automatically transfer to a stronger range
topology.  Put
\[
 E_2=\ell^2,\qquad
 E_1=\{y:\sum n^4|y_n|^2<\infty\},\qquad
 A x=(n^{-2}x_n)_n.
\]
The inclusion $T:E_1\to E_2$ is continuous, $A$ is trace class on
$E_2$, and its range is exactly $E_1$.  Its unique lift
$A_0:E_2\to E_1$ is an isometry onto $E_1$, hence is not compact
and is not nuclear.  This verifies the logical distinction needed
in the application on p.~515.

For the finite, smooth, annular real-place operator studied above,
the missing control \emph{can} be provided: its Schwartz kernel is
flat in the output coordinate near zero.  The Hermite expansion
in the input coordinate has coefficient functions in
$\mathcal S_{00}$ rapidly decreasing in every Schwartz seminorm.
Writing the expansion with scalar weights $(1+n)^{-2}$ gives a
nuclear map $L^2\to\mathcal S_{00}$ (the rescaled coefficient
functions are bounded in that Fr\'echet space and the Hermite
coefficient functionals have norm one on $L^2$).  Thus the local
regularization has a genuine nuclear lift.  This is a proof for
the explicitly specified real operator.  It does not verify an
unspecified lift from the global weighted Hilbert spaces, nor
establish the global-to-$S$-local operator identity required in
the missing Section~12.3.

\paragraph{Remaining gate.}
The proved local repairs are the centered prime weight and the
annular rate $\lambda'/\lambda''\to0$, with sufficiently thin
boundary smoothing.  An arithmetic transport still has to carry
the global prime contribution into the local comparison and to
control the relevant Hilbert or nuclear trace norms.  None of
MA-LA-001--006 claims that this is impossible.  Conversely, a
correct local cutoff limit alone does not provide that transport.
